\documentclass[10pt,a4paper]{article}
\usepackage{fontspec}
\usepackage{polyglossia}
\setmainlanguage{russian}
\setotherlanguage{english}
\setmainfont{Liberation Serif}
\usepackage[left=1.0cm,right=1.0cm,top=0.8cm,bottom=0.7cm]{geometry}
\usepackage{amsmath,amssymb}
\setlength{\parindent}{0pt}
\pagestyle{empty}

% ШПАРГАЛКА ПРИНИМАЮЩЕМУ. Одна страница, трём преподавателям на стол.
% ОТСТУПЛЕНИЕ ОТ КРИТЕРИЯ ХЭНДОФФА, осознанное: там сказано «только с ответами».
% Здесь есть ещё короткая идея решения и строка «что спрашивать» — потому что
% листок принимается ПО РЕШЕНИЮ, и принимающему нужно за десять секунд понять,
% верной ли дорогой идёт школьник. Каждая идея умещается в две-три строки.

\newcommand{\zz}[3]{\par\vspace{5pt}\noindent%
  \textbf{#1.}\ \textbf{#2}\newline\small #3\normalsize\par}

\begin{document}

\noindent{\large\bfseries Шпаргалка принимающему}\hfill{\small Кружок, 7\,И\quad·\quad23 сентября\quad·\quad листок «Один ответ, несколько, ни одного»}
\vspace{3pt}\hrule\vspace{6pt}

\noindent\small
Листок принимается \textbf{по решению}. Ответ в поле — для самопроверки школьника, а не основание для плюса.
Порядок задач — по возрастанию. Три линии: \textbf{1\,→\,6} утверждения о самих себе; \textbf{3\,→\,4} конфигурация с параметром; \textbf{2\,→\,5\,→\,7} доказать, что иначе нельзя.
Буква $n$ стоит только в задачах 4 и 6 — там от неё зависит тип ответа. В задаче 7 числа нет вовсе: чётность там не дана, а доказывается.\normalsize

\vspace{6pt}\hrule

\zz{1}{Два ответа: либо ложны все десять, либо истинно только последнее.}
{Любые два утверждения противоречат друг другу $\Rightarrow$ истинных не больше одного. Ноль истинных: ложных десять, а утверждения говорят про 0…9 ложных — ни одно не попало, всё сходится. Одно истинное: ложных девять, это последнее.
\textbf{Что спрашивать:} нашёл ли оба случая. Один случай — не решение.}

\zz{2}{Ни одного рыцаря, все 30 — лжецы.}
{Возьмём самого левого рыцаря: слева от него рыцарей ноль, справа $\ge 0$, значит его фраза ложна — противоречие. Проверка: при нуле рыцарей у всех «$0>0$» ложно, все честно врут.
\textbf{Что спрашивать:} сделал ли проверку. Доказать, что рыцарей нет, — половина дела.}

\zz{3}{67.}
{Трёх рыцарей подряд нет: средний утверждал бы, что соседи разных типов, а они оба рыцари. Значит в любых трёх подряд не больше двух. Разбиваем 100 на 33 тройки и одного: $33\cdot2+1=67$. Пример: Р\,ЛРР\,ЛРР\,…\,ЛРР.
\textbf{Что спрашивать:} есть ли обе половины. Пример без оценки и оценка без примера — каждое половина; «наибольшее» без доказательства не доказано.}

\zz{4}{(а) При любых $n$. \quad (б) Если $n$ делится на 3 — ноль или $2n/3$; иначе только ноль.}
{У рыцаря ровно один сосед-рыцарь $\Rightarrow$ рыцари стоят блоками ровно по два. Два лжеца рядом невозможны: у крайнего оказался бы ровно один сосед-рыцарь. Значит по кругу ЛРР, период 3. Отдельно: все лжецы подходят всегда.
\textbf{Типичная потеря:} забывают вариант «все лжецы» и в (а), и в (б).}

\zz{5}{Рыцарей больше.}
{Считаем пары «рыцарь\,--\,лжец» двумя способами. Рыцарь даёт ровно одну $\Rightarrow$ пар ровно $R$. Лжец солгал, значит дружит хотя бы с одним рыцарем $\Rightarrow$ пар $\ge L$. Отсюда $R\ge L$; жителей нечётно, значит $R\ne L$, то есть $R>L$.
\textbf{Что спрашивать:} где использована нечётность. Без неё получается только $R\ge L$.}

\zz{6}{При чётных $n$ ложных ровно $n/2$. При нечётных $n$ такого быть не может.}
{Если $k$-е истинно, истинны все предыдущие $\Rightarrow$ истинные образуют начало. Пусть их $m$, ложных $n-m$. Из истинности $m$-го: $n-m\ge m$. Из ложности $(m{+}1)$-го: $n-m\le m$. Значит $n=2m$.
\textbf{Что спрашивать:} проверены ли края $m=0$ и $m=n$. Оба невозможны.}

\zz{7$^\star$}{Доказательство. Попутно выходит, что жителей чётное число.}
{Проведём стрелку от каждого к тому, на кого он указал. Из каждого выходит ровно одна стрелка и в каждого входит ровно одна $\Rightarrow$ идя по стрелкам, обязательно вернёмся туда, откуда вышли; все жители разбиваются на такие замкнутые обходы. Вдоль обхода типы чередуются: рыцарь указывает на лжеца, лжец — на рыцаря. Чередование замыкается только при чётной длине обхода, значит в каждом обходе рыцарей и лжецов поровну. Складываем по всем обходам.
\textbf{Что спрашивать:} почему, идя по стрелкам, мы вернёмся в начало. Слова «цикл» семиклассник не скажет — принимать «вернёмся туда, откуда вышли».
\textbf{Проверено перебором:} при нечётном числе жителей расстановки не существует вовсе; при чётном во всех расстановках рыцарей ровно половина.}

\zz{8$^\star$}{Идём с конца.}
{(а) Если бы дожили до 31-го, накануне все бы знали — значит не 31-го. (б) Раз 31-е отпало, последний возможный день 27-е, и рассуждение повторяется. (в) Так же отпадает 20-е.
(г) Парадокс: рассуждение (а)–(в) опирается на то, что объявлению верят; поверив ему, кружковцы исключают все дни — и именно поэтому 20-го оказываются застигнуты врасплох. Формального ответа нет, и это честный итог.
\textbf{Как считать:} (а), (б), (в) — по пункту каждый. (г) не ждём.}

\vspace{6pt}\hrule\vspace{4pt}

\noindent\small\textbf{ЗАПАС.} В файле листка закомментированы две задачи средней сложности, обе проверены перебором.
\textbf{Первая:} 25 островитян, каждый сказал «число рыцарей среди нас чётно» — \emph{ответ: так не бывает} (чётное $k$ $\Rightarrow$ фразу произнесли и лжецы; нечётное $k$ $\Rightarrow$ её произнесли и рыцари; оба случая невозможны).
\textbf{Вторая:} 100 человек, чай любят 70, кофе 75, сок 80, молоко 85 — \emph{хотя бы 10 любят всё} (не любящих не больше $30{+}25{+}20{+}15=90$).
Ставить их — только взамен двух из середины: иначе коридор сумм уедет за пределы 3–7.

\vspace{3pt}
\noindent\small\textbf{ЗАПАСНОЙ ГРОБ}, если седьмая окажется лёгкой. Дурак считает всех дураками, а себя умным; скромный умный про всех знает верно, а себя считает дураком; уверенный умный знает верно про всех и про себя. В думе 200 депутатов; премьер спросил каждого из 199 присутствующих, сколько умных в зале, и по анкетам ответа не понял. Вернулся двухсотый, заполнил анкету про всю думу включая себя — и премьер понял. Сколько умных в думе? — \emph{ответ: 1 или 3}. Проверено перебором для залов от 3 до 7 человек: ответ от размера зала не зависит.\normalsize

\end{document}
